\chapter{绪论与数学基础}

\intro{
	数学是上帝用来书写宇宙的语言。
	
	\rightline{——Galileo Galilei}
}

\section{系统动力学建模概述}

\subsection{系统动力学的基本概念}

系统动力学建模与仿真，是以系统的观点，将工程对象视为由若干相互关联的元件组成的整体，通过建立数学模型来描述其动态行为，并利用数值计算方法进行模拟求解的一门学科。

\begin{myTheorem}{系统动力学建模的一般步骤}
	\begin{enumerate}
		\item \textbf{物理建模}：分析系统的力学特性，识别关键变量与参数；
		\item \textbf{数学建模}：根据物理定律建立描述系统行为的微分方程；
		\item \textbf{数值求解}：选择适当的数值方法对微分方程进行离散求解；
		\item \textbf{仿真分析}：编写程序代码，对系统在不同条件下的响应进行模拟；
		\item \textbf{验证与优化}：将仿真结果与理论解或实验数据对比，修正模型。
	\end{enumerate}
\end{myTheorem}

\subsection{动力学系统的分类}

\begin{mydefinition}{动力学系统}
	动力学系统是指状态随时间演化的系统，其数学描述通常为微分方程或微分代数方程组。
\end{mydefinition}

按不同标准，动力学系统可分类如下：
\begin{enumerate}
	\item \textbf{按时间变量}：连续时间系统（常微分方程）与离散时间系统（差分方程）；
	\item \textbf{按线性性质}：线性系统与非线性系统；
	\item \textbf{按自由度}：单自由度系统（SDOF）与多自由度系统（MDOF）；
	\item \textbf{按参数特性}：时不变系统（LTI）与时变系统（LTV）；
	\item \textbf{按激励方式}：自由振动系统与受迫振动系统。
\end{enumerate}

\section{常微分方程基本概念}

\subsection{常微分方程的定义与分类}

\begin{mydefinition}{常微分方程}
	含有未知函数及未知函数导数的方程称为微分方程。若方程中的未知函数仅依赖于一个自变量，则称为\textbf{常微分方程}（ODE）。
\end{mydefinition}

$n$ 阶常微分方程的一般形式为：
\begin{myformula}
	F\left(t, y, y', y'', \ldots, y^{(n)}\right) = 0
\end{myformula}

若能解出最高阶导数，可写为显式形式：
\begin{myformula}
	y^{(n)} = f\left(t, y, y', \ldots, y^{(n-1)}\right)
\end{myformula}

\begin{mydefinition}{线性与非线性}
	若 $F$ 关于 $y, y', \ldots, y^{(n)}$ 为线性，则方程为\textbf{线性常微分方程}；否则为\textbf{非线性常微分方程}。
\end{mydefinition}

线性 $n$ 阶常微分方程的标准形式为：
\begin{myformula}
	a_n(t) y^{(n)} + a_{n-1}(t) y^{(n-1)} + \cdots + a_1(t) y' + a_0(t) y = g(t)
\end{myformula}

\begin{remark}
	线性方程具有叠加原理：若 $y_1, y_2$ 为齐次方程的解，则它们的线性组合 $c_1 y_1 + c_2 y_2$ 也是齐次方程的解。
\end{remark}

\subsection{初值问题与边值问题}

\begin{mydefinition}{初值问题（IVP）}
	给定 $n$ 阶微分方程及在初始点 $t_0$ 处的 $n$ 个初始条件：
	\begin{myformula}
		\begin{cases}
			y^{(n)} = f(t, y, y', \ldots, y^{(n-1)}) \\[4pt]
			y(t_0) = y_0, \; y'(t_0) = y_0', \; \ldots, \; y^{(n-1)}(t_0) = y_0^{(n-1)}
		\end{cases}
	\end{myformula}
\end{mydefinition}

一阶常微分方程组的初值问题可统一写为：
\begin{myformula}
	\frac{d\mathbf{y}}{dt} = \mathbf{f}(t, \mathbf{y}), \quad \mathbf{y}(t_0) = \mathbf{y}_0
\end{myformula}

\section{解的存在唯一性}

\begin{theorem}{Picard 存在唯一性定理}
	设函数 $f(t, y)$ 在矩形区域 $R = \{(t, y) \mid |t - t_0| \le a, |y - y_0| \le b\}$ 内连续，且关于 $y$ 满足\textbf{Lipschitz 条件}：
	\[
	|f(t, y_1) - f(t, y_2)| \le L |y_1 - y_2|, \quad \forall (t, y_1), (t, y_2) \in R
	\]
	则初值问题 $y' = f(t, y), \; y(t_0) = y_0$ 在区间 $|t - t_0| \le \min(a, b/M)$ 上存在唯一的解，其中 $M = \max_{(t,y)\in R} |f(t, y)|$。
\end{theorem}

\begin{proof}
	\textbf{Step 1：构造 Picard 迭代序列。}
	定义迭代序列 $\{y_n(t)\}$：
	\[
	y_0(t) \equiv y_0, \quad y_{n+1}(t) = y_0 + \int_{t_0}^{t} f(s, y_n(s)) \, ds
	\]
	
	\textbf{Step 2：证明序列的一致收敛性。}
	利用 Lipschitz 条件，可证：
	\[
	|y_{n+1}(t) - y_n(t)| \le M L^n \frac{|t - t_0|^{n+1}}{(n+1)!}
	\]
	由 Weierstrass M-判别法，级数 $\sum [y_{n+1}(t) - y_n(t)]$ 一致收敛，故 $\{y_n(t)\}$ 一致收敛于某个极限函数 $y(t)$。
	
	\textbf{Step 3：证明极限函数满足积分方程。}
	由一致收敛性，可对迭代式取极限：
	\[
	y(t) = y_0 + \int_{t_0}^{t} f(s, y(s)) \, ds
	\]
	对两端求导即得 $y'(t) = f(t, y(t))$，且 $y(t_0) = y_0$。
	
	\textbf{Step 4：证明唯一性。}
	设有两个解 $y(t)$ 和 $\tilde{y}(t)$，则：
	\[
	|y(t) - \tilde{y}(t)| \le \int_{t_0}^{t} |f(s, y(s)) - f(s, \tilde{y}(s))| \, ds \le L \int_{t_0}^{t} |y(s) - \tilde{y}(s)| \, ds
	\]
	由 Gronwall 不等式，得 $|y(t) - \tilde{y}(t)| \equiv 0$，故解唯一。
\end{proof}

\begin{corollary}{Lipschitz 条件的充分条件}
	若 $f(t, y)$ 在 $R$ 上具有连续偏导数 $\frac{\partial f}{\partial y}$，则 $f$ 在 $R$ 上关于 $y$ 满足 Lipschitz 条件。
\end{corollary}

\begin{proof}
	由微分中值定理：
	\[
	f(t, y_1) - f(t, y_2) = \frac{\partial f}{\partial y}(t, \xi) (y_1 - y_2)
	\]
	取 $L = \max_{(t,y)\in R} \left|\frac{\partial f}{\partial y}\right|$，即得 Lipschitz 条件。
\end{proof}

\subsection{Gronwall 不等式}

\begin{theorem}{Gronwall 不等式（微分形式）}
	设 $\alpha(t), u(t)$ 为 $[a,b]$ 上的连续函数，且 $\alpha(t) \ge 0$。若
	\[
	u(t) \le C + \int_{a}^{t} \alpha(s) u(s) \, ds, \quad \forall t \in [a,b]
	\]
	则
	\[
	u(t) \le C \, \exp\left(\int_{a}^{t} \alpha(s) \, ds\right)
	\]
\end{theorem}

\begin{proof}
	令 $v(t) = C + \int_{a}^{t} \alpha(s) u(s) \, ds$，则 $v(a) = C$ 且 $v'(t) = \alpha(t) u(t)$。由条件 $u(t) \le v(t)$ 及 $\alpha(t) \ge 0$：
	\[
	v'(t) = \alpha(t) u(t) \le \alpha(t) v(t)
	\]
	即 $v'(t) - \alpha(t) v(t) \le 0$。乘以积分因子 $\exp(-\int_{a}^{t} \alpha(s) ds)$：
	\[
	\frac{d}{dt}\left[ v(t) e^{-\int_{a}^{t} \alpha(s) ds} \right] \le 0
	\]
	积分得 $v(t) e^{-\int_{a}^{t} \alpha(s) ds} \le v(a) = C$，故 $u(t) \le v(t) \le C e^{\int_{a}^{t} \alpha(s) ds}$。
\end{proof}

\section{将高阶方程化为一阶方程组}

\begin{proposition}{降阶方法}
	任意 $n$ 阶显式微分方程
	\[
	y^{(n)} = f(t, y, y', \ldots, y^{(n-1)})
	\]
	均可等价地化为 $n$ 个一阶方程组成的方程组。
\end{proposition}

\begin{proof}
	引入新变量：
	\[
	z_1 = y, \quad z_2 = y', \quad z_3 = y'', \quad \ldots, \quad z_n = y^{(n-1)}
	\]
	则原方程等价于：
	\[
	\begin{cases}
		z_1' = z_2 \\
		z_2' = z_3 \\
		\;\; \vdots \\
		z_{n-1}' = z_n \\
		z_n' = f(t, z_1, z_2, \ldots, z_n)
	\end{cases}
	\]
	即 $\mathbf{z}' = \mathbf{F}(t, \mathbf{z})$，其中 $\mathbf{z} = (z_1, \ldots, z_n)^T$。
\end{proof}

\begin{example}{弹簧-质量-阻尼系统}
	单自由度弹簧-质量-阻尼系统的运动方程为：
	\[
	m \ddot{x} + c \dot{x} + k x = F(t)
	\]
	令 $x_1 = x$（位移），$x_2 = \dot{x}$（速度），化为一阶方程组：
	\[
	\begin{cases}
		\dot{x}_1 = x_2 \\[4pt]
		\dot{x}_2 = -\dfrac{k}{m} x_1 - \dfrac{c}{m} x_2 + \dfrac{F(t)}{m}
	\end{cases}
	\]
	此即系统的\textbf{状态空间表示}。
\end{example}

\begin{remark}
	降阶的意义在于：只需研究一阶方程组的数值解法，即可处理任意高阶微分方程。这极大简化了数值算法的设计。
\end{remark}
